\scanpage{109}

a) compatibilité au chgt de base

b) \deleted{\unsure{$f(-\omega,x,y)=f(y,x)$}} \hfill $f_{-\omega}(x,y)=f_\omega(y,x)$

c) $f_\omega(\deleted{x}y,\deleted{y}x)=\sigma f_\omega(\deleted{y}x,\deleted{x}y)$

d) $f_\omega(x,x)=i(x)$

\margin{donc b) et c) ensemble se lisent $f_{-\omega}(x,y)=f_\omega(y,x)=\sigma f_\omega(\unsure{x,y})$}

\deleted{\unsure{Si on} \unsure{Nous}}

Notons que

\[
\delta'(-q)=-\delta'(q)=-\omega^{\otimes2}
\]

donc \deleted{il y a un iso} on peut trouver, grâce à $\omega$, une algèbre de \deleted{\unsure{Azumaya}}\unsure{quaternions} $\mathcal M$ et un iso

\[
E\overset{\unsure{\varphi|\psi}}{\simeq}\unsure{\gamma|\mathfrak g}\mathcal M
\]

transformant \uuline{$-q$} en $\unsure{\cdot}q_{\mathcal M}=-\det\mid\gamma\mathcal M$ et $\omega$ en $\omega_{\mathcal M}$, donc transformant $q$ en $\det$ ; \inserted{$(\mathcal M,\varphi)$ est défini à isom unique près par $q,\omega$}

% slanted

\margin{Attention au signe}

Ici $X$ s'interprète comme la variété projective associée au cône des él. \inserted{\unsure{str$^{\mathrm t}$}} nilpotents de $\mathcal M$, définis par les équations

\[
\operatorname{Tr}u=\det u=0
\]

Quant à $\Sigma$, il correspond à l'équation $\det u=-1$ en plus de $\operatorname{Tr}u=0$, qui correspond au schém. des \uline{réflexions} associées à $\mathcal M$. L'iso. induit par $f_\omega$ doit être un iso

\[
(X\times X)\setminus X\simeq\Sigma
\]

s'interprétant comme une bijection (fonctorielle

% archive sheet number: 108
